\documentclass[a4paper,12pt]{article}
\usepackage[utf8]{inputenc}
\usepackage[spanish]{babel}
\usepackage{geometry}
\usepackage{amsmath}
\usepackage{graphicx}
\usepackage{circuitikz}
\usetikzlibrary{babel}
\usepackage{enumitem}
\usepackage{float}
\usepackage{caption}
\usepackage{subcaption}

\geometry{top=2.5cm, bottom=2.5cm, left=2.5cm, right=2.5cm}

\begin{document}

\begin{center}
    \textbf{PRACTICA 26} \\[0.5cm]
    \textbf{\large TEOREMA DE THEVENIN}
\end{center}

\vspace{0.5cm}

\noindent \textbf{FINALIDADES}

\begin{enumerate}
    \item Determinar analíticamente los valores de $V_{\text{TH}}$ (tensión del generador de Thévenin) y $R_{\text{TH}}$ (resistencia del generador de Thévenin) en un circuito de c.c. que contiene una sola fuente de tensión.
    \item Confirmar experimentalmente los valores de $V_{\text{TH}}$ y $R_{\text{TH}}$ propuestos por el teorema de Thévenin en la solución de circuitos puente desequilibrados.
\end{enumerate}

\vspace{0.5cm}

\noindent \textbf{INFORMACION PRELIMINAR}

Hemos empleado las leyes de Ohm y Kirchhoff para resolver los problemas que plantean los circuitos de c.c. En el estudio de los circuitos eléctricos y electrónicos complicados son necesarias a veces otras bases de cálculo. El teorema de Thévenin forma parte de un grupo de teoremas de redes (incluyendo el de superposición y el de Norton) que proporciona los medios adecuados para simplificar el análisis de circuitos complicados. La técnica empleada implica reducir una red a un circuito sencillo equivalente que actúa como la red original.

\vspace{0.5cm}

\noindent \textbf{Teorema de Thévenin}

El teorema de Thévenin enuncia que cualquier red de dos terminales se puede substituir por un circuito sencillo equivalente consistente en un generador de Thévenin, cuya tensión $V_{\text{TH}}$, actuando en serie con una resistencia interna $R_{\text{TH}}$, hace que pase corriente por la carga. La figura 26-1 \textit{d} es el equivalente de Thévenin para el circuito de la figura 26-1 \textit{a} considerando a $R_{\text{L}}$ como carga. Si conocemos los valores $V_{\text{TH}}$ y $R_{\text{TH}}$ el proceso de hallar la corriente $I$ en $R_{\text{L}}$ se reducirá a una sencilla aplicación de la ley de Ohm.

\vspace{0.5cm}

\begin{figure}[H]
    \centering
    % Row 1: (a) and (b)
    \begin{subfigure}[b]{0.45\textwidth}
        \centering
        \begin{circuitikz}[scale=0.8, transform shape]
            % Source and Series resistors
            \draw (0,3) to[battery1, l=100 V] (0,0);
            \draw (0,3) to[R, l=$R_1\ 5\,\Omega$] (3,3); % Top resistor
            \draw (0,0) to[R, l=$R_2\ 195\,\Omega$] (3,0); % Bottom resistor
            
            % Parallel Resistor R3
            \draw (3,3) to[R, l=$R_3\ 200\,\Omega$, *-*] (3,0);
            
            % Terminals A and B labels
            \node[above] at (3,3) {A};
            \node[below] at (3,0) {B};
            
            % Load resistor RL
            \draw (3,3) -- (5,3) to[R, l=$R_L\ 350\,\Omega$] (5,0) -- (3,0);
            
            \node at (2.5,-0.8) {(a)};
        \end{circuitikz}
    \end{subfigure}
    \hfill
    \begin{subfigure}[b]{0.45\textwidth}
        \centering
        \begin{circuitikz}[scale=0.8, transform shape]
            % Source and Series resistors
            \draw (0,3) to[battery1, l=100 V] (0,0);
            \draw (0,3) to[R, l=$R_1\ 5\,\Omega$] (3,3); 
            \draw (0,0) to[R, l=$R_2\ 195\,\Omega$] (3,0); 
            
            % Parallel Resistor R3
            \draw (3,3) to[R, l=$R_3\ 200\,\Omega$, *-*] (3,0);
            
            % Open terminals A and B
            \draw (3,3) -- (4,3);
            \draw (3,0) -- (4,0);
            \node[above] at (4,3) {A};
            \node[below] at (4,0) {B};
            
            % Vth measurement
            \draw[<->] (4.5,3) -- (4.5,0) node[midway, right] {$V_{\text{TH}} = 50\,\text{V}$}; % Corrected closing brace
            
            \node at (2.5,-0.8) {(b)};
        \end{circuitikz}
    \end{subfigure}
    
    \vspace{1cm}
    
    % Row 2: (c) and (d)
    \begin{subfigure}[b]{0.45\textwidth}
        \centering
        \begin{circuitikz}[scale=0.8, transform shape]
            % Shorted Source
            \draw (0,3) -- (0,0); % Short
            \draw (0,3) to[R, l=$R_1\ 5\,\Omega$] (3,3); 
            \draw (0,0) to[R, l=$R_2\ 195\,\Omega$] (3,0); 
            
            % Parallel Resistor R3
            \draw (3,3) to[R, l=$R_3\ 200\,\Omega$, *-*] (3,0);
            
            % Looking in Logic (arrow)
            \draw (3,1.5) to[open, v={}] (5,1.5);
            \draw[thick, <-] (4.5,1.5) -- (3.5,1.5) node[above, pos=1.2] {Vista};
            % Terminals
            \node[above] at (3,3) {A};
            \node[below] at (3,0) {B};
            
            % Equation text below or to side, let's put it as text in figure or label
            % The image has equation: Rth = 200(195+5)/(200+(195+5)) = 100
             \node[align=left, right, scale=0.8] at (0,-1.5) {$R_{\text{TH}} = \frac{200(195+5)}{200+(195+5)} = 100\,\Omega$};
            
            \node at (2.5,-2.5) {(c)};
        \end{circuitikz}
    \end{subfigure}
    \hfill
    \begin{subfigure}[b]{0.45\textwidth}
        \centering
        \begin{circuitikz}[scale=0.8, transform shape]
            % Thevenin Equivalent
            % Battery Vth and Rth in series
            \draw (0,3) to[battery1, l=$V_{\text{TH}}\ 50\,\text{V}$] (0,0);
            \draw (0,3) to[R, l=$R_{\text{TH}}\ 100\,\Omega$] (3,3);
            
            % Terminals A and B
            \node[above] at (3,3) {A};
            \node[below] at (3,0) {B};
            \draw (3,3) to[short, *-] (3,3); % dot
            \draw (3,0) -- (0,0); % return path
            \draw (3,0) to[short, *-] (3,0); % dot
            
            % Load
            \draw (3,3) -- (5,3) to[R, l=$R_{\text{L}}\ 350\,\Omega$] (5,0) -- (3,0);
            
            % Current I
            \draw[->] (3.5, 3.2) -- (4.5, 3.2) node[midway, above] {$I$};
            
            % Equation
            \node[align=left, below, scale=0.8] at (3.5,-0.5) {$I = \frac{V_{\text{TH}}}{R_{\text{TH}} + R_{\text{L}}} = \frac{50}{450} = 0{,}111\,\text{A}$};
            
            \node at (2.5,-1.8) {(d)};
        \end{circuitikz}
    \end{subfigure}
    
    \caption{Análisis de una red por el teorema de Thévenin.}
    \label{fig:26-1};
\end{figure}

\end{document}
